% Probability / expected value - the formulas worth having on paper
\section{Probability \& Expected Value}

\subsection{Basic Rules}
\[P(\overline{A})=1-P(A)\qquad
  P(\text{at least one})=1-P(\text{none})\]
\[P(\ge 1\text{ success in }n\text{ trials})=1-(1-p)^n\]
\[P(A\cup B)=P(A)+P(B)-P(A\cap B)\]
\[P(A\mid B)=\frac{P(A\cap B)}{P(B)}\qquad
  P(A\cap B)=P(A)\,P(B\mid A)\]
\noindent Only when independent: $P(A\cap B)=P(A)P(B)$.

\subsection{Bayes / Total Probability}
\noindent For a partition $A_1,\dots,A_n$ of the sample space:
\[P(B)=\sum_i P(B\mid A_i)\,P(A_i)\]
\[P(A_i\mid B)=\frac{P(B\mid A_i)\,P(A_i)}{\sum_j P(B\mid A_j)\,P(A_j)}\]

\subsection{Distributions}
\noindent\textbf{Binomial} --- $k$ successes in $n$ trials:
\[P(X=k)=\binom{n}{k}p^k(1-p)^{n-k}\]
\[E[X]=np\qquad \operatorname{Var}(X)=np(1-p)\]
\noindent\textbf{Geometric} --- first success on trial $k$:
\[P(X=k)=(1-p)^{k-1}p\qquad E[X]=\frac{1}{p}\]
\noindent Cost $c$ per attempt $\Rightarrow E[\text{cost}]=c/p$.

\noindent\textbf{Negative binomial} --- $r$-th success on trial $n$:
\[P(X=n)=\binom{n-1}{r-1}p^r(1-p)^{n-r}\qquad E[X]=\frac{r}{p}\]

\subsection{Expected Value}
\[E[X]=\sum_x x\,P(X=x)\qquad E[aX+b]=aE[X]+b\]
\[E[f(X)]=\sum_x f(x)\,P(X=x)\qquad E[X^2]\neq\bigl(E[X]\bigr)^2\]
\noindent\textbf{Linearity} --- holds even when the $X_i$ are dependent:
\[E\Bigl[\sum_i X_i\Bigr]=\sum_i E[X_i]\]

\subsection{Indicators}
\noindent $I_A=1$ if $A$ happens, else $0$, so $E[I_A]=P(A)$. For
$X=\sum_i I_i$:
\[E[X]=\sum_i P(A_i)\]
\noindent Expected number of valid pairs:
\[E[X]=\sum_{i<j}P\bigl((i,j)\text{ valid}\bigr)\]
\noindent Expected distinct values over $n$ draws:
\[E[\text{distinct}]=\sum_v\bigl[1-(1-p_v)^n\bigr]\]
\noindent Asked for an expected count? Reach for indicators plus
linearity before building the whole distribution.

\subsection{Variance}
\[\operatorname{Var}(X)=E[X^2]-\bigl(E[X]\bigr)^2\]
\[\operatorname{Var}(X+Y)=\operatorname{Var}(X)+\operatorname{Var}(Y)
  +2\operatorname{Cov}(X,Y)\]
\noindent Independent $\Rightarrow
\operatorname{Var}\bigl(\sum_i X_i\bigr)=\sum_i\operatorname{Var}(X_i)$.
Variance needs that independence; linearity of $E$ never does.

\subsection{Expected-Value DP}
\noindent From state $s$, move to $t$ with probability $p_t$ paying
$c(s,t)$:
\[E[s]=\sum_t p_t\bigl(c(s,t)+E[t]\bigr)\]
\noindent One operation per step:
\[E_i=1+\sum_j p_{ij}E_j\]
\noindent Self-loop --- finish with probability $p$, else repeat, cost
$c$ each time:
\[E=pc+(1-p)(c+E)\ \Longrightarrow\ E=\frac{c}{p}\]
\noindent Many states $\Rightarrow$ solve $A\mathbf{x}=\mathbf{b}$ by
Gaussian elimination.
